\section{Fairness Definitions}

\begin{frame}{Turning ``Fair'' into Something You Can Measure}

    \begin{block}{Notation for this section}
        \begin{itemize}
            \item $A$ --- the group attribute (e.g.\ $A \in \{\GrpA, \GrpB\}$).
            \item $Y \in \{0,1\}$ --- the true outcome (repaid the loan, re-offended, was ill).
            \item $R \in [0,1]$ --- the score the model outputs.
            \item $\Yhat = \mathbf{1}[R > t]$ --- the decision after thresholding.
        \end{itemize}
    \end{block}

    \vspace{0.8em}

    Almost every group-fairness definition is one of three statements:

    \vspace{0.4em}
    \centering
    \small
    \begin{tabular}{lll}
        \toprule
        \textbf{Family} & \textbf{Statement} & \textbf{Plain reading} \\
        \midrule
        Independence & $\Yhat \perp A$ & the decision rate is the same for every group \\
        Separation   & $\Yhat \perp A \mid Y$ & the error rates are the same for every group \\
        Sufficiency  & $Y \perp A \mid R$ & a score means the same thing for every group \\
        \bottomrule
    \end{tabular}

    \vspace{1em}
    \raggedright
    {\scriptsize Framing after S.~Barocas, M.~Hardt and A.~Narayanan, \emph{Fairness and Machine
    Learning: Limitations and Opportunities}, MIT Press, 2023,
    \href{https://fairmlbook.org/}{fairmlbook.org}, Chapter 3.}
\end{frame}
